\begin{figure}[ht]
\centering
\begin{tikzpicture}[x=1cm, y=1cm, font=\small]

% Curve illustrating local min, local max, inflection, endpoint optimum
% Use f(x) = 0.05 x (x-3)(x-7)(x-11) + 1.5, shifted/scaled into [0, 12] x [0, 4.5]
% Actually use a hand-tuned cubic-like polynomial showing two interior crit pts + endpoints

% Axes
\draw[->, thick] (-0.3, 0) -- (12.5, 0) node[right] {$x$};
\draw[->, thick] (0, -0.2) -- (0, 5.0) node[above] {$f(x)$};

% Domain markers
\draw[engblack, dashed, thin] (1, 0) -- (1, 5.0);
\draw[engblack, dashed, thin] (11, 0) -- (11, 5.0);
\node[engblack, anchor=north, font=\scriptsize] at (1, 0) {$x = a$};
\node[engblack, anchor=north, font=\scriptsize] at (11, 0) {$x = b$};

% A smooth curve: f(x) = 1.5 + sin((x - 1)*pi/5) + 0.4*sin((x-1)*pi/2.5)
\draw[very thick, engblue]
  plot[domain=1:11, samples=120]
    (\x, {2.2 + 1.0*sin((\x - 1)*36) + 0.4*sin((\x - 1)*72)});

% Compute features by inspection at sampled points
% At x = 3.4 (local max), x = 6 (inflection w/ horizontal tangent? not really, just inflection), x = 8.6 (local min)
% Just place markers approximately
\fill[engred] (3.4, {2.2 + 1.0*sin(2.4*36) + 0.4*sin(2.4*72)}) circle (2.5pt);
\node[engred, anchor=south, font=\scriptsize, align=center]
  at (3.4, {2.2 + 1.0*sin(2.4*36) + 0.4*sin(2.4*72) + 0.05})
  {local max\\ $f'=0,\,f''<0$};

\fill[engamber] (8.6, {2.2 + 1.0*sin(7.6*36) + 0.4*sin(7.6*72)}) circle (2.5pt);
\node[engamber, anchor=north, font=\scriptsize, align=center]
  at (8.6, {2.2 + 1.0*sin(7.6*36) + 0.4*sin(7.6*72) - 0.1})
  {local min\\ $f'=0,\,f''>0$};

% Endpoint markers
\fill[engblack] (1, {2.2}) circle (2pt);
\node[engblack, anchor=south east, font=\scriptsize, align=right]
  at (0.9, {2.2 + 0.05}) {left\\ endpoint};

\fill[engblack] (11, {2.2 + 1.0*sin(10*36) + 0.4*sin(10*72)}) circle (2pt);
\node[engblack, anchor=south west, font=\scriptsize, align=left]
  at (11.1, {2.2 + 1.0*sin(10*36) + 0.4*sin(10*72) + 0.05}) {right\\ endpoint};

% Horizontal tangents at the critical points
\draw[engred, thick] (2.6, {2.2 + 1.0*sin(2.4*36) + 0.4*sin(2.4*72)})
                  -- (4.2, {2.2 + 1.0*sin(2.4*36) + 0.4*sin(2.4*72)});
\draw[engamber, thick] (7.8, {2.2 + 1.0*sin(7.6*36) + 0.4*sin(7.6*72)})
                   -- (9.4, {2.2 + 1.0*sin(7.6*36) + 0.4*sin(7.6*72)});

% Caption legend
\node[anchor=north, font=\scriptsize, align=left]
  at (6, -0.7)
  {On $[a, b]$, the extrema lie at interior critical points (Fermat: $f'(x) = 0$, classify by $f''$) or at the endpoints.\\
   The exhaustive procedure: list all candidates, evaluate $f$ at each, pick the largest and smallest.};

\end{tikzpicture}
\caption[Critical points, second-derivative classification, and
endpoint candidates]{One-variable optimisation on a closed interval
$[a, b]$. Interior local extrema satisfy Fermat's necessary condition
$f'(x) = 0$, with the second-derivative test sorting maxima ($f'' < 0$)
from minima ($f'' > 0$). Endpoints are candidates regardless of the
derivative. The global maximum and minimum lie at one of the
candidates; the procedure evaluates $f$ at each and picks the
extremes. Points where $f'(x) = 0$ but $f''(x) = 0$ require higher
derivatives to classify, and points where $f$ is not differentiable
(kinks, cusps) must be added to the candidate list directly.}
\label{fig:vol02:ch05:optimisation-extrema}
\end{figure}
